%!TEX root = ../../main.tex
\section{선행 연구의 위치}
\label{sec:rel-position}

저자의 선행 연구 \cite{baek2026killconditioned}는 공개 텔레메트리에서 교전 전 정보의 회복 가능성을 물었고, 교전 승자를 교환가치 라벨로 정의하여 공학적 테이블 요약 위의 그래디언트 부스팅이 신경망 표현들보다 우수함을 보고하였다. 그 연구는 킬 수만을 예측한 것이 아니며, 고정 계수 교환가치와 운영적 라벨의 한계를 스스로 명시하였고, 같은 테이블 입력의 MLP 비교와 시간순 패치 분할도 이미 포함하였다. 본 논문은 이를 처음 도입했다고 주장하지 않는다.

본 논문이 선행 연구와 달라지는 지점은 세 가지이다. 첫째, 정의의 근거: 손으로 정한 상수 대신 자료에서 추정한 경계와 규칙에 고정한 게이트, 그리고 상수마다의 출처 분류이다. 둘째, 결과의 정의: 교전 국소적 교환가치 대신 경기 승패와 연결된 동결 평가기의 변화 방향이다. 셋째, 예측의 기준: 절대 성능이 아니라 승률·시간의 유연한 기준선을 넘는 이득이다. 라벨과 코호트가 달라졌으므로 선행 연구의 AUC와 본 논문의 SVI AUC의 차이를 개선 효과라고 부르지 않는다. 선행 연구의 교전 승자 예측 계보에서 사후에 도출되었던 ``규모가 클수록 예측 가능하다''는 결론은 시간 특징의 누출로 철회되었고(제 \ref{sec:eng-scale} 절), 본 논문은 그 철회를 유지한다.
